% !TeX root = ../main.tex
\section{Mechanism Study and Pre-Registered Field Study: Sheaf Energy versus
Contradiction Counting}
\label{sec:mechanism}

The framework's mathematical identity rests on one empirical question: does
the sheaf energy $f_3$ add signal beyond the trivial baseline that counts
NLI-flagged contradiction edges above a decision threshold? This section
answers it in two stages: a simulation that isolates the \emph{mechanism} by
which a separation could arise (\S\ref{sec:simulation}),
and the pre-registered \emph{field study} that executed the protocol with a
real, calibrated instrument on real reasoning chains (\S\ref{sec:field}). The
field study adjudicates the simulation's assumptions, confirms
the pre-registered success criteria where the separation is structural, and
sharpens the claim into Prop.~\ref{prop:struct-equiv}.

\subsection{Mechanism simulation}
\label{sec:simulation}

The simulation isolates the mechanism under two documented assumptions about
NLI behavior:

\begin{itemize}
\item[(S1)] \emph{Distance decay.} The confidence assigned by a pairwise NLI
model to the direct contradiction edge decays with the semantic/reasoning
distance $k$ between the violating statement and the claim it contradicts; in
the simulation, $c(k)=\mathrm{clip}(0.95-0.13k+\mathcal N(0,0.06))$, falling
below the standard $0.5$ decision threshold around $k\ge 3$.
\item[(S2)] \emph{Noise in both classes.} Hedged or hypothetical claims
(considered and rejected by the text, hence not asserted) are flagged by NLI
as contradicting asserted claims with confidence $w\sim U(0.55,0.85)$; and
spurious sub-threshold asserted--asserted contradiction edges occur with
$w\sim U(0.10,0.45)$.
\end{itemize}

\paragraph{Setup.} Synthetic reasoning chains of six asserted claims connected
by entailment edges ($w\sim\mathrm{Beta}(20,2)$); violated cases append an
asserted statement contradicting claim $A_0$ at step distance
$k\in\{1,\dots,5\}$, with the direct edge confidence drawn per (S1) and an
optional weak paraphrase edge to $A_1$; both classes receive (S2) noise. 120
cases per $(k,\text{class})$ cell. Three detectors: \textsc{Count}
(\#contradiction edges with $w\ge 0.5$; the standard baseline), \textsc{Sheaf}
(the unified energy of \S\ref{sec:logic-gate}, assertion-aware, no
thresholding), \textsc{Sat} (DPLL on edges with $w\ge 0.5$).

\begin{table}[t]
\centering
\begin{tabular}{cccccc}
\toprule
$k$ & AUROC \textsc{Count} & AUROC \textsc{Sheaf} & AUROC \textsc{Sat} &
Loc.\ \textsc{Count} & Loc.\ \textsc{Sheaf}\\
\midrule
1 & 0.881 & 1.000 & 1.000 & 100.0\% & 100.0\%\\
2 & 0.875 & 1.000 & 1.000 & 100.0\% & 100.0\%\\
3 & 0.831 & 1.000 & 0.933 & 86.7\% & 100.0\%\\
4 & 0.537 & 0.993 & 0.554 & 10.8\% & 97.5\%\\
5 & 0.504 & 0.970 & 0.500 & 0.0\% & 83.3\%\\
\bottomrule
\end{tabular}
\caption{Mechanism simulation: detection AUROC and top-1 localization rate
versus violation distance $k$ ($n=120$ per cell), \emph{under simulated
confidences drawn from (S1)/(S2)}. Threshold-based methods collapse to chance
as the direct-edge confidence crosses the threshold; the weighted sheaf energy
degrades gracefully and retains localization.}
\label{tab:simulation}
\end{table}

The simulation establishes a conditional statement: \emph{if} real NLI
confidences decay with distance (S1) and produce hedged-claim false positives
(S2), \emph{then} the sheaf energy separates structurally from counting
(Table~\ref{tab:simulation}). Both assumptions are directly testable, and the
harness becomes the real experiment by replacing simulated confidences with
model outputs. That experiment follows.

\subsection{Pre-registered field study}
\label{sec:field}

\paragraph{Design (protocol frozen before data collection; deviations
logged).} Instrument: DeBERTa-large \cite{he2021} fine-tuned on MultiNLI
\cite{williams2018}, temperature-scaled \cite{guo2017} on MNLI
validation-matched ($n=2{,}000$; $T=1.452$; 15-bin ECE $0.011$ post- vs.\
$0.047$ pre-calibration; MNLI accuracy $0.917$); the instrument never receives
injection locations. Corpus: 450 GSM8K \cite{cobbe2021} gold solutions parsed
into 4--10-step claim chains, with the dataset's native calculator annotations
as the numeric extraction channel. Paired clean/violated design, $n=150$ cases per cell (45
for per-channel threshold calibration, 105 for evaluation); 22{,}706 scored
sentence pairs. Violation families: \textbf{V1/V2}, template negation of the
step $k\in\{1,\dots,5\}$ positions before the end, appended at the end;
\textbf{V3}, a compositional numeric violation---a closing sentence misstates
the ratio between final and initial computed quantities while every sentence
\emph{pair} remains individually consistent (pairwise detectors blind by
construction); \textbf{V4}, hedged contamination attached as a non-asserted
vertex in both classes. Arms: \textsc{Count}, \textsc{Wsum} (de-thresholded
weighted sum---the control that prevents attributing a sheaf win to mere
de-thresholding), \textsc{NliMax}, \textsc{Sat}, and the max-form sheaf energy
of Def.~\ref{def:maxform}. Primary endpoints: paired-bootstrap
$\Delta$AUROC against the best baseline per cell with Holm correction across
cells; top-1 localization. Pre-registered success criteria: $\Delta$AUROC
$\ge 0.03$ (CI excluding 0) or localization advantage $\ge 0.10$, plus a
\textsc{Wsum}-specific win on V3 or V4.

\paragraph{Assumption adjudication (E0).} Both simulation assumptions
\emph{fail} in the field, in informative ways. (S1): calibrated contradiction
confidence on the direct pair $(\neg c_i,c_i)$ is \emph{flat} in $k$
($0.749,0.740,0.742,0.764,0.756$ for $k=1,\dots,5$)---a modern NLI recognizes
a verbatim negation regardless of its position; cross-pair confidence sits
near $0.5$ with a mild downward drift. (S2): hedged constructions were flagged
in $0/300$ trials (Clopper--Pearson 95\%: $[0,0.012]$). Consequently the
collapse mechanism of Table~\ref{tab:simulation} is not exercised by verbatim
negations, and hedged reconciliation has nothing to discount under this
instrument.

\begin{table}[t]
\centering
\small
\begin{tabular}{lcccccc c}
\toprule
Cell & \textsc{Sheaf}$_{\max}$ & \textsc{Sheaf}$_{\mathrm{quad}}$ &
\textsc{Count} & \textsc{Wsum} & \textsc{NliMax} & \textsc{Sat} &
Loc.\ S/B\\
\midrule
V1 $k{=}1$ & 0.719 & 0.719 & 0.696 & 0.719 & 0.531 & 0.514 & 0.21 / 0.94\\
V2 $k{=}2$ & 0.741 & 0.741 & 0.710 & 0.741 & 0.520 & 0.505 & 0.17 / 0.91\\
V2 $k{=}3$ & 0.756 & 0.756 & 0.734 & 0.756 & 0.513 & 0.505 & 0.06 / 0.94\\
V2 $k{=}4$ & 0.717 & 0.717 & 0.685 & 0.717 & 0.503 & 0.500 & 0.03 / 0.90\\
V2 $k{=}5$ & 0.739 & 0.739 & 0.701 & 0.739 & 0.502 & 0.500 & 0.02 / 0.92\\
\textbf{V3} & \textbf{0.990} & 0.586 & 0.498 & 0.501 & 0.500 & 0.500 &
\textbf{0.24 / 0.00}\\
V4 $k{=}1$ & 0.751 & 0.751 & 0.696 & 0.732 & 0.518 & 0.510 & 0.14 / 0.91\\
\bottomrule
\end{tabular}
\caption{Field study (confirmatory run; $n=105$ evaluation cases per cell).
AUROC per arm; Loc.\ S/B is top-1 localization for the sheaf hotspot vs.\ the
baseline above-threshold rule. On V3 the sheaf--best-baseline gap is
$\Delta=+0.489$ [95\% CI $0.470,0.503$], Holm-adjusted $p=0.0018$, with
FPR at 95\% TPR of $1.0\%$ (vs.\ $94.3\%$ for every pairwise baseline); the
\textsc{Wsum} anti-artifact clause holds. V1/V2 equalities are exact
(Prop.~\ref{prop:struct-equiv}).}
\label{tab:field}
\end{table}

\paragraph{Findings (Table~\ref{tab:field}).} Three results delimit the
tool's value precisely.

\emph{(i) The structural regime: confirmed at $0.990$.} On compositional
violations (V3), every pairwise arm is at chance by construction---no sentence
pair witnesses the inconsistency---while the max-form sheaf energy detects the
broken cycle holonomy at AUROC $0.990$ with $1\%$ FPR at $95\%$ TPR, and is
the only method with nonzero localization ($0.24$ vs.\ exactly $0$). The
pre-registered criteria are met with the \textsc{Wsum} clause, so the
separation is not a de-thresholding artifact.

\emph{(ii) The structure-free regime: provable coincidence.} On V1/V2, the
sheaf and \textsc{Wsum} columns agree to three decimals in every cell---an
instance of Prop.~\ref{prop:struct-equiv}, not an empirical accident. Field
NLI noise (spurious contradiction edges between ordinary solution steps) caps
all pairwise methods at ${\approx}0.72$--$0.76$ and drives \textsc{Sat} to
chance; sheaf top-1 localization loses to the simple above-threshold rule
outside V3 for the same reason. The localization claim is therefore scoped to
compositional violations.

\emph{(iii) The aggregation form is forced.} The as-specified quadrature
combination of the two $f_3$ channels scores $0.586$ on V3; the max-form of
Def.~\ref{def:maxform} scores $0.990$ from the same channel values. The
dilution phenomenon of Prop.~\ref{prop:dilution} recurs inside $f_3$, and its
P3 remedy applies one level down.

\subsection{Field test of the calibration layer}
\label{sec:field-e2}

On clean-chain sheaf scores (pool $n=735$; $n_{\mathrm{cal}}=200$,
$\alpha=0.10$, $R=1{,}000$ random splits), pooled exceedance is
$0.0952\le\alpha$ (exact binomial $p=1.0$ against inflation): the
finite-sample bound of Theorem~\ref{thm:split} holds in the field (and again
in the replication and scale-up runs: $0.0995$ and $0.0996$). A naive KS test
of per-split coverages against the $\mathrm{Beta}(m,n+1-m)$ refinement
rejects ($p\approx 0$), but null checks show this reflects the audit
statistic, not the law: the same KS rejects ideal continuous i.i.d.\ scores
under the same procedure (test-set binomial noise, pool reuse, and, here,
$39.7\%$ tied scores), while the analytically exact conditional coverage
passes (KS $p=0.64$). See \S\ref{sec:calibration} for the corrected audit
prescription; the tie-safe exceedance bound is the field-verified object.

\paragraph{Scale-up.} A doubled-scale run ($1{,}000$ chains, $n=210$
evaluation cases per cell, six negation and three closing-sentence template
families, $46{,}936$ scored pairs) reproduces every finding: V3 at AUROC
$0.995$ ($\Delta=+0.488$ [$0.474,0.498$], FPR at 95\% TPR $0.5\%$,
localization $0.30$ vs.\ $0.00$), exact V1/V2 counting-coincidence, quadrature
dilution ($0.576$), and the Regime-I bound ($0.0996\le 0.10$, pool
$n=1{,}470$). The headline is therefore stable across sample size and
template families.

\paragraph{Replication with a second instrument.} The key cells were
replicated with an architecture-distinct instrument
(RoBERTa-large \cite{liu2019} fine-tuned on MNLI, $T=1.337$, ECE $0.012$;
600 chains; $n=150$/cell). The
headline replicates: on V3, the max-form sheaf energy scores AUROC
$\mathbf{0.999}$ (FPR at 95\% TPR $0.5\%$; localization $0.33$ vs.\ $0.00$)
against chance-level baselines ($\Delta=+0.489$ [$0.479,0.498$], Holm
$p=0.0008$); the quadrature dilution ($0.593$) and the structural coincidence
on V1 (max-form $\equiv$ \textsc{Wsum} $=0.752$) also replicate; the
Regime-I bound holds again ($0.0995\le 0.10$). Instructively, this
instrument's hedge false-positive rate is \emph{nonzero} ($2.3\%$,
Clopper--Pearson $[0.9\%,4.8\%]$), and exactly as the mechanism predicts, the
assertion-aware reconciliation capability activates: on V4 the sheaf beats
\textsc{Wsum} by $+0.035$ ($\ge$ the pre-registered $0.03$ bar) where under
the zero-S2 instrument it was vacuous. The V4 capability is thus an
\emph{instrument-conditional} one, switched on precisely when
$\mathrm{S2}>0$---a falsifiable dependence the instruments jointly
demonstrate as a dose--response pattern: $\mathrm{S2}=0.7\%$ (DeBERTa,
scale-up) yields $+0.015$ over \textsc{Wsum}, $\mathrm{S2}=2.3\%$ (RoBERTa)
yields $+0.035$.

\paragraph{Scope of the claim.} Violations are template-inserted into real
chains (the pre-registered method); an organically generated error corpus is
registered as the immediate next step. All code, the frozen protocol with its
deviation log, and per-run artifacts are released with the paper.
